% ---------------------------------------------------------------- sections/07-conclusion.tex
\section{Conclusion}\label{sec:concl}

Our verdict: the corporate income tax makes financing matter, but it does not make debt a free lunch. In the two-state firm, debt of 40 moves 0.73 of value from the tax authority to investors, which vanishes without the tax (attribution), and no investor can copy it with homemade leverage, because only the firm can deduct interest. The money machine stops at the debt the weak year can repay, $\Dstar = 44.4$, and as long as distress destroys more than about 3.6\% of the cash available in default, borrowing past that limit lowers firm value (robustness); estimated distress costs of 10--20\% \citep{andrade1998} are well above this threshold. What we did not expect is that a higher tax rate makes each unit of debt more valuable but lowers the safe amount of debt, because the tax also comes out of the weak state's cash. In our presentation we called the prize small, about 1--2\% of firm value; that is the one-period result, and once debt is maintained the same mechanism is worth about 10\% (Section~\ref{sec:multi}), close to \citeauthor{graham2000}'s \citeyearpar{graham2000} estimate. Unlike a trial of the equal-borrowing-terms assumption, ours shows homemade leverage failing even when investors borrow on the same terms as firms. Table~\ref{tab:summary} shows that each extension and application changes the size or location of the gain; none makes debt free, and only investor-level taxes can remove the tax advantage of interest altogether.

Our analysis has limitations. Valuation uses risk-neutral probabilities. In the two-state model this affects only the price of distress, but in Section~\ref{sec:calib} we treat historical EBIT distributions as if they were risk-neutral. Distress tends to occur when the whole market is weak, so its risk-neutral probability exceeds its historical frequency \citep{almeida2007}; we therefore understate the present value of distress costs and overstate optimal debt. In addition, EBIT is drawn independently each year; distress costs are a fixed proportion of value; and the firm holds no cash buffer, so a single bad year triggers default. We also abstract from other determinants of capital structure, including agency costs, asymmetric information and the option value of unused debt capacity, which Section~\ref{sec:calib} suggests are most important for firms at either extreme. Our conclusion would be weakened by evidence that firms with larger interest deductions are not worth more, by distress costs below about 3.6\% of default-state cash, or by a marginal investor who is taxed far more heavily on interest than on equity income. For practitioners, the relevant question is less ``what is my tax rate?'' than ``what is the worst year through which the firm must still service its debt?''

\begin{table}[htbp]\centering\footnotesize
\caption{Summary of results}\label{tab:summary}
\begin{tabularx}{\linewidth}{@{}>{\raggedright\arraybackslash\hsize=0.64\hsize}X
                                >{\raggedright\arraybackslash\hsize=1.54\hsize}X
                                >{\centering\arraybackslash\hsize=0.82\hsize}X@{}}
\toprule
Analysis & Main finding & Verdict holds? \\
\midrule
Continuous cash flows (Section~\ref{sec:cont}) & Interior optimum at $D \approx 34$ with 1\% default probability; gain 0.85\% of $\VU$ & Yes \\
\addlinespace[1.5pt]
Parameter sensitivity (Section~\ref{sec:sens}) & Debt limit between 37 and 52; gain 1--2\% of $\VU$; debt past the limit destroys value if distress costs exceed about 3.6\% of default-state cash & Yes; location of the limit is fragile \\
\addlinespace[1.5pt]
Market-priced debt (Section~\ref{sec:mkt}) & Optimum almost unchanged; shareholders rather than lenders bear expected distress costs & Yes \\
\addlinespace[1.5pt]
Usability of deductions (Section~\ref{sec:usable}) & Effective shield below $\tauc$ for volatile or highly levered firms; a 30\% cap binds below coverage of 3.3 in an average year & Yes; adds a second limit \\
\addlinespace[1.5pt]
Personal taxes (Section~\ref{sec:personal}) & Advantage per unit of interest from $-1.7\%$ (US top bracket) to $\tauc$ (Singapore, tax-exempt investors) & Depends on the investor \\
\addlinespace[1.5pt]
Multi-period debt (Section~\ref{sec:multi}) & Gain rises to about 10\% of value; optimal leverage falls with EBIT volatility & Yes; gain is larger \\
\addlinespace[1.5pt]
US TCJA (Section~\ref{sec:tcja}) & Rate cut lowers the value that debt adds by about 40\%; an EBIT-based cap reduces optimal debt by up to a third & Yes \\
\addlinespace[1.5pt]
Singapore and S-REITs (Section~\ref{sec:sg}) & Shield reaches investors in full; S-REITs approximate the no-tax benchmark under a regulatory coverage floor & Yes \\
\addlinespace[1.5pt]
Calibration to two firms (Section~\ref{sec:calib}) & Volatile airline: no tax case for debt, yet S\$12.9bn borrowed; stable supermarket: about 12\% of value forgone & Ranking yes; actual choices differ; SIA's ``no debt'' corner reflects our no-buffer default rule \\
\bottomrule
\end{tabularx}
\tablenote{\textit{Notes:} ``Verdict'' means that corporate tax makes debt valuable up to a limit set by the firm's worst state, and never free.}
\end{table}
